% GENERATED FILE. Edit canonical lessons, then run npm run generate:books.
% canonical-source-hash: fnv1a64:a5607c52715c3f10
% canonical-lessons: ML-C351-sampresanam, ML-C351-srotavu, ML-C351-ripporttar, ML-C351-citrikarikkuka, ML-C351-sampresanam-ceyyuka

\chapter{On the Air}
\label{ch:ml-on-the-air}
\hlchaptermodality{351}

\begin{chapteropening}
\textbf{By the end of this chapter:} \emph{I can say a broadcast, a listener, a reporter, to depict and to broadcast.}

Complete the last lesson of chapter 351: I can say a broadcast, a listener, a reporter, to depict and to broadcast.
\end{chapteropening}

\section[\texorpdfstring{\ml{സംപ്രേഷണം} (samprēṣaṇaṁ)}{സംപ്രേഷണം (samprēṣaṇaṁ)}]{\ml{സംപ്രേഷണം} (samprēṣaṇaṁ) — a broadcast; a transmission}

\label{lesson:ML-C351-sampresanam}



\begin{quote}
Before the new one: say the Malayalam for a reader, then the Malayalam for to print.
\end{quote}

\subsection*{You'll want to know: \ml{സംപ്രേഷണം}}
\textbf{\ml{സംപ്രേഷണം}} — \emph{samprēṣaṇaṁ} — \textquotedblleft{}a broadcast; a transmission\textquotedblright{}.

\begin{quote}
\ml{മത്സരത്തിന്റെ} \ml{സംപ്രേഷണം} \ml{രാത്രി} \ml{എട്ട്} \ml{മണിക്ക്} \ml{തുടങ്ങും}.
\end{quote}

\begin{quote}
\emph{matsarattinṟe samprēṣaṇaṁ rātri eṭṭŭ maṇikkŭ tuṭaṅṅuṁ.}
\end{quote}

\begin{quote}
The broadcast of the match will begin at eight at night.
\end{quote}

A broadcast goes out on a \ml{ചാനൽ}, a channel, at a set time.

\begin{cousinweb}
\textbf{\ml{സംപ്രേഷണം}} (\emph{samprēṣaṇaṁ}) is borrowed from Sanskrit \emph{saṁpreṣaṇa}, \textquotedblleft{}the sending forth or dispatching of something\textquotedblright{} (Monier-Williams), from \emph{pra-iṣ}, \textquotedblleft{}to send off\textquotedblright{}. A broadcast is a programme sent out to everyone at once. It is often written \ml{സംപ്രേക്ഷണം}, as if it came from \ml{പ്രേക്ഷകൻ}, a viewer, but the Sanskrit word behind it is the one for sending.
\end{cousinweb}

\begin{grammarlens}[title={What you've built}]
One more word for telling what happened.
\end{grammarlens}

\subsection*{Your turn}
\begin{itemize}
\raggedright
  \item \emph{Say it:} \emph{samprēṣaṇaṁ}
  \item \emph{Say it:} \emph{samprēṣaṇaṁ}, once more
  \item \emph{Say it:} \emph{matsarattinṟe samprēṣaṇaṁ rātri eṭṭŭ maṇikkŭ tuṭaṅṅuṁ.} — The broadcast of the match will begin at eight at night.
  \item \emph{Recall:} say the Malayalam for a reader, then the Malayalam for to print, then say \emph{samprēṣaṇaṁ} again
\end{itemize}

\begin{tcolorbox}[breakable,colback=teal!4,colframe=teal!35!black,title={Before you move on}]
What does \ml{സംപ്രേഷണം} mean? (\textquotedblleft{}A broadcast; a transmission\textquotedblright{}.) Say it once more.
\end{tcolorbox}
\section[\texorpdfstring{\ml{ശ്രോതാവ്} (śrōtāvŭ)}{ശ്രോതാവ് (śrōtāvŭ)}]{\ml{ശ്രോതാവ്} (śrōtāvŭ) — a listener}

\label{lesson:ML-C351-srotavu}



\begin{quote}
Before the new one: say the Malayalam for to print, then the Malayalam for a broadcast.
\end{quote}

\subsection*{You'll want to know: \ml{ശ്രോതാവ്}}
\textbf{\ml{ശ്രോതാവ്}} — \emph{śrōtāvŭ} — \textquotedblleft{}a listener\textquotedblright{}.

\begin{quote}
\ml{ആ} \ml{പരിപാടിക്ക്} \ml{ഗ്രാമത്തിൽ} \ml{ഒരുപാട്} \ml{ശ്രോതാക്കൾ} \ml{ഉണ്ട്}.
\end{quote}

\begin{quote}
\emph{ā paripāṭikkŭ grāmattil orupāṭŭ śrōtākkaḷ uṇṭŭ.}
\end{quote}

\begin{quote}
That programme has many listeners in the village.
\end{quote}

Plural: \ml{ശ്രോതാക്കൾ}, \emph{śrōtākkaḷ}.

Plural \ml{ശ്രോതാക്കൾ}. A viewer watches; a listener hears.

\begin{cousinweb}
\textbf{\ml{ശ്രോതാവ്}} (\emph{śrōtāvŭ}) is borrowed from Sanskrit \emph{śrotṛ}, \textquotedblleft{}one who hears, a hearer\textquotedblright{} (Monier-Williams). Sanskrit nouns in -tṛ come into Malayalam in -\ml{താവ്}.
\end{cousinweb}

\begin{grammarlens}[title={What you've built}]
One more word for telling what happened.
\end{grammarlens}

\subsection*{Your turn}
\begin{itemize}
\raggedright
  \item \emph{Say it:} \emph{śrōtāvŭ}
  \item \emph{Say it:} \emph{śrōtāvŭ}, once more
  \item \emph{Say it:} \emph{ā paripāṭikkŭ grāmattil orupāṭŭ śrōtākkaḷ uṇṭŭ.} — That programme has many listeners in the village.
  \item \emph{Recall:} say the Malayalam for to print, then the Malayalam for a broadcast, then say \emph{śrōtāvŭ} again
\end{itemize}

\begin{tcolorbox}[breakable,colback=teal!4,colframe=teal!35!black,title={Before you move on}]
What does \ml{ശ്രോതാവ്} mean? (\textquotedblleft{}A listener\textquotedblright{}.) Say it once more.
\end{tcolorbox}
\section[\texorpdfstring{\ml{റിപ്പോർട്ടർ} (ṟippōrṭṭar)}{റിപ്പോർട്ടർ (ṟippōrṭṭar)}]{\ml{റിപ്പോർട്ടർ} (ṟippōrṭṭar) — a reporter}

\label{lesson:ML-C351-ripporttar}



\begin{quote}
Before the new one: say the Malayalam for a broadcast, then the Malayalam for a listener.
\end{quote}

\subsection*{You'll want to know: \ml{റിപ്പോർട്ടർ}}
\textbf{\ml{റിപ്പോർട്ടർ}} — \emph{ṟippōrṭṭar} — \textquotedblleft{}a reporter\textquotedblright{}.

\begin{quote}
\ml{റിപ്പോർട്ടർ} \ml{ക്യാമറയുമായി} \ml{ഗ്രാമത്തിൽ} \ml{എത്തി}.
\end{quote}

\begin{quote}
\emph{ṟippōrṭṭar kyāmaṟayumāyi grāmattil etti.}
\end{quote}

\begin{quote}
The reporter reached the village with a camera.
\end{quote}

The spoken word beside the formal \ml{ലേഖകൻ}.

\begin{cousinweb}
\textbf{\ml{റിപ്പോർട്ടർ}} (\emph{ṟippōrṭṭar}) was borrowed from English \emph{reporter}. Malayalam had already taken \ml{റിപ്പോർട്ട്} (\emph{ṟippōrṭṭŭ}), \textquotedblleft{}a report\textquotedblright{}, from English; the person who writes one followed.
\end{cousinweb}

\begin{grammarlens}[title={What you've built}]
One more word for telling what happened.
\end{grammarlens}

\subsection*{Your turn}
\begin{itemize}
\raggedright
  \item \emph{Say it:} \emph{ṟippōrṭṭar}
  \item \emph{Say it:} \emph{ṟippōrṭṭar}, once more
  \item \emph{Say it:} \emph{ṟippōrṭṭar kyāmaṟayumāyi grāmattil etti.} — The reporter reached the village with a camera.
  \item \emph{Recall:} say the Malayalam for a broadcast, then the Malayalam for a listener, then say \emph{ṟippōrṭṭar} again
\end{itemize}

\begin{tcolorbox}[breakable,colback=teal!4,colframe=teal!35!black,title={Before you move on}]
What does \ml{റിപ്പോർട്ടർ} mean? (\textquotedblleft{}A reporter\textquotedblright{}.) Say it once more.
\end{tcolorbox}
\section[\texorpdfstring{\ml{ചിത്രീകരിക്കുക} (citrīkarikkuka)}{ചിത്രീകരിക്കുക (citrīkarikkuka)}]{\ml{ചിത്രീകരിക്കുക} (citrīkarikkuka) — to depict; to film}

\label{lesson:ML-C351-citrikarikkuka}



\begin{quote}
Before the new one: say the Malayalam for a listener, then the Malayalam for a reporter.
\end{quote}

\subsection*{You'll want to know: \ml{ചിത്രീകരിക്കുക}}
\textbf{\ml{ചിത്രീകരിക്കുക}} — \emph{citrīkarikkuka} — \textquotedblleft{}to depict; to film\textquotedblright{}.

\begin{quote}
\ml{അവർ} \ml{ആ} \ml{സിനിമ} \ml{ഒരു} \ml{ചെറിയ} \ml{ഗ്രാമത്തിൽ} \ml{ചിത്രീകരിച്ചു}.
\end{quote}

\begin{quote}
\emph{avar ā sinima oru ceṟiya grāmattil citrīkariccu.}
\end{quote}

\begin{quote}
They filmed that film in a small village.
\end{quote}

Past: \ml{ചിത്രീകരിച്ചു}, \emph{citrīkariccu}.

Built on \ml{ചിത്രം}, a picture: to make into a picture, on paper or on film.

\begin{cousinweb}
\textbf{\ml{ചിത്രീകരിക്കുക}} (\emph{citrīkarikkuka}) is built from Sanskrit \emph{citrīkaraṇa}, \textquotedblleft{}making bright, painting\textquotedblright{} (Monier-Williams), from \emph{citra}, \textquotedblleft{}a picture\textquotedblright{}, with the Malayalam verb ending -\ml{ിക്കുക}.
\end{cousinweb}

\begin{grammarlens}[title={What you've built}]
One more word for telling what happened.
\end{grammarlens}

\subsection*{Your turn}
\begin{itemize}
\raggedright
  \item \emph{Say it:} \emph{citrīkarikkuka}
  \item \emph{Say it:} \emph{citrīkarikkuka}, once more
  \item \emph{Say it:} \emph{avar ā sinima oru ceṟiya grāmattil citrīkariccu.} — They filmed that film in a small village.
  \item \emph{Recall:} say the Malayalam for a listener, then the Malayalam for a reporter, then say \emph{citrīkarikkuka} again
\end{itemize}

\begin{tcolorbox}[breakable,colback=teal!4,colframe=teal!35!black,title={Before you move on}]
What does \ml{ചിത്രീകരിക്കുക} mean? (\textquotedblleft{}To depict; to film\textquotedblright{}.) Say it once more.
\end{tcolorbox}
\section[\texorpdfstring{\ml{സംപ്രേഷണം} \ml{ചെയ്യുക} (samprēṣaṇaṁ ceyyuka)}{സംപ്രേഷണം ചെയ്യുക (samprēṣaṇaṁ ceyyuka)}]{\ml{സംപ്രേഷണം} \ml{ചെയ്യുക} (samprēṣaṇaṁ ceyyuka) — to broadcast; to air}

\label{lesson:ML-C351-sampresanam-ceyyuka}



\begin{quote}
Before the new one: say the Malayalam for a reporter, then the Malayalam for to depict.
\end{quote}

\subsection*{You'll want to know: \ml{സംപ്രേഷണം} \ml{ചെയ്യുക}}
\textbf{\ml{സംപ്രേഷണം} \ml{ചെയ്യുക}} — \emph{samprēṣaṇaṁ ceyyuka} — \textquotedblleft{}to broadcast; to air\textquotedblright{}.

\begin{quote}
\ml{ചാനൽ} \ml{ആ} \ml{മത്സരം} \ml{രാത്രി} \ml{സംപ്രേഷണം} \ml{ചെയ്തു}.
\end{quote}

\begin{quote}
\emph{cānal ā matsaraṁ rātri samprēṣaṇaṁ ceytu.}
\end{quote}

\begin{quote}
The channel broadcast the match at night.
\end{quote}

Past: \ml{സംപ്രേഷണം} \ml{ചെയ്തു}, \emph{samprēṣaṇaṁ ceytu}.

A verb in its own right, worth learning whole: the noun with \ml{ചെയ്യുക}, to do, as so many formal verbs are made.

\begin{cousinweb}
\textbf{\ml{സംപ്രേഷണം} \ml{ചെയ്യുക}} (\emph{samprēṣaṇaṁ ceyyuka}) is built from \ml{സംപ്രേഷണം} (\emph{samprēṣaṇaṁ}), \textquotedblleft{}a broadcast\textquotedblright{}, borrowed from Sanskrit \emph{saṁpreṣaṇa}, \textquotedblleft{}a sending forth\textquotedblright{} (Monier-Williams), and \ml{ചെയ്യുക}, \textquotedblleft{}to do\textquotedblright{}.
\end{cousinweb}

\begin{grammarlens}[title={What you've built}]
One more word for telling what happened.
\end{grammarlens}

\subsection*{Your turn}
\begin{itemize}
\raggedright
  \item \emph{Say it:} \emph{samprēṣaṇaṁ ceyyuka}
  \item \emph{Say it:} \emph{samprēṣaṇaṁ ceyyuka}, once more
  \item \emph{Say it:} \emph{cānal ā matsaraṁ rātri samprēṣaṇaṁ ceytu.} — The channel broadcast the match at night.
  \item \emph{Recall:} say the Malayalam for a reporter, then the Malayalam for to depict, then say \emph{samprēṣaṇaṁ ceyyuka} again
\end{itemize}

\begin{tcolorbox}[breakable,colback=teal!4,colframe=teal!35!black,title={Before you move on}]
What does \ml{സംപ്രേഷണം} \ml{ചെയ്യുക} mean? (\textquotedblleft{}To broadcast; to air\textquotedblright{}.) Say it once more.
\end{tcolorbox}
